% tnde user guide -- for physics users.
% Build:  pdflatex user_guide.tex && pdflatex user_guide.tex     (twice, for the TOC)
\documentclass[11pt,a4paper]{article}

\usepackage[T1]{fontenc}
\usepackage[utf8]{inputenc}
\usepackage{lmodern}
\usepackage[margin=1in]{geometry}
\usepackage{amsmath,amssymb}
\usepackage{booktabs}
\usepackage{array}
\usepackage{tabularx}
\usepackage{ragged2e}
\usepackage{xcolor}
\usepackage{tcolorbox}
\usepackage{enumitem}
\usepackage{parskip}
\usepackage{listings}
\usepackage[hidelinks]{hyperref}

\definecolor{codebg}{gray}{0.96}
\definecolor{rule}{gray}{0.55}

\lstset{
  language=Python,
  basicstyle=\ttfamily\small,
  backgroundcolor=\color{codebg},
  frame=none,
  columns=fullflexible,
  keepspaces=true,
  xleftmargin=1em,
  framexleftmargin=1em,
  framexrightmargin=1em,
  framextopmargin=4pt,
  framexbottommargin=4pt,
  showstringspaces=false,
  breaklines=true,
  keywordstyle=\ttfamily,
  commentstyle=\color{rule},
  stringstyle=\ttfamily,
}

% Inline code.  Underscores are escaped by hand at every call site.
\newcommand{\code}[1]{\texttt{#1}}

\newcommand{\kcut}{k_{\mathrm{cut}}}
\newcommand{\kmax}{k_{\max}}
\newcommand{\dd}{\mathrm{d}}
\newcommand{\vx}{\mathbf{x}}
\newcommand{\vk}{\mathbf{k}}

\newtcolorbox{finding}{
  colback=codebg, colframe=rule, boxrule=0.5pt, arc=1pt,
  left=8pt, right=8pt, top=6pt, bottom=6pt,
}

\newcolumntype{L}{>{\RaggedRight\arraybackslash}X}

\title{\textbf{\texttt{tnde} user guide}\\[2pt]
  \large Time-dependent PDEs on quantics tensor trains, for physics users}
\date{}

\begin{document}
\maketitle
\thispagestyle{empty}

\noindent
This guide is for someone who has a Schr\"odinger, Gross--Pitaevskii, diffusion or
reaction--diffusion problem and wants to run it on a grid that would not fit in memory.
It explains what the code computes, what the tensor-network representation does to the
field, which knobs matter physically and which are numerical, and how to get numbers
back out that can be trusted. The \code{README} covers installation and the validation
tables; \code{docs/guide.tex} covers the original Gross--Pitaevskii benchmark in
depth. This document sits between them.

\vspace{1em}
\hrule
\vspace{0.5em}
{\small\tableofcontents}
\vspace{0.5em}
\hrule
\newpage

% ---------------------------------------------------------------------------
\section{What the code solves}

Every equation \code{tnde} handles has the form
\begin{equation}
\partial_t u(\vx,t)
= \underbrace{\ell(-i\nabla)\,u}_{\text{linear, translation-invariant}}
+ \underbrace{p(\vx)\,u}_{\text{local linear}}
+ \underbrace{N(u,\vx,t)}_{\text{local nonlinear}}
\label{eq:master}
\end{equation}
on a periodic box in one, two or three space dimensions, for a scalar or a
multi-component field $u$. The three terms are the three things the user supplies.

\begin{tabularx}{\textwidth}{@{}lLL@{}}
\toprule
term & meaning & examples \\
\midrule
$\ell(\vk)$ & the \emph{symbol} of a linear operator diagonal in momentum: in Fourier
space $\partial_t \hat u_{\vk} = \ell(\vk)\,\hat u_{\vk}$ &
$-i\hbar k^2/2m$ (free Schr\"odinger), $-Dk^2$ (diffusion), $-\nu k^4$
(hyperviscosity), $-D|k|^{\alpha}$ (fractional diffusion, L\'evy flights),
$-i\,\mathbf c\cdot\vk$ (advection) \\
$p(\vx)$ & a multiplicative term that depends only on position &
$-iV(\vx)/\hbar$ (a potential in real time), $-V(\vx)$ (in imaginary time), a
spatially varying growth rate $r(\vx)$ \\
$N(u,\vx,t)$ & at each point, an ODE in $u$ alone &
$-ig|u|^2u$ (mean field), $ru(1-u)$ (logistic growth), $u-u^3$ (double well), any
coupling between components \\
\bottomrule
\end{tabularx}

The one restriction is that $N$ contains \textbf{no derivatives}. Burgers ($u\,u_x$),
Korteweg--de Vries and Navier--Stokes are outside the scheme, because the splitting of
Section~\ref{sec:step} needs the nonlinear part to act pointwise. Everything else that
can be written as \eqref{eq:master} is one
\code{Equation(linear=..., potential=..., nonlinear=...)} away, and the common cases
are presets.

The point of the tensor-train representation is that the cost is set by how much
\emph{structure} the solution has, not by how many grid points there are. A 1D
condensate on $2^{30}\approx10^9$ points, or a 2D field on $4096^2$ points, is evolved
without the grid ever existing in memory. Section~\ref{sec:quantics} says exactly which
fields this is good for.

% ---------------------------------------------------------------------------
\section{A first run}

Install once:
\begin{lstlisting}[language=bash]
pip install numpy scipy jax
\end{lstlisting}
That is the whole list. The tensor-train and cross-interpolation kernels come from
\code{qutecipy}, which is vendored inside the repository, so a plain clone runs anywhere
with nothing else installed. Put the repository directory on \code{sys.path} and both
packages resolve (the examples do this themselves).

The ground state of a trapped 1D Bose gas by imaginary-time evolution, on a grid of
$2^{20}$ points:
\begin{lstlisting}
import numpy as np
from tnde import Grid, equations, pde

grid = Grid(20, (-16.0, 16.0))                         # 2**20 points on [-16, 16)
eq = equations.gross_pitaevskii(V=lambda x: 0.5 * x**2, g=50.0, imaginary_time=True)
psi, snaps = pde.solve(eq, lambda x: np.exp(-x**2 / 8), grid,
                       dt=0.01, nsteps=400, tolerance=1e-10, maxdim=30, save_every=100)

(x,), values = pde.reconstruct(psi, grid, n=1024)      # the field on 1024 sample points
density = np.abs(values)**2
n0 = np.abs(pde.evaluate(psi, grid, 0.0))**2           # the density at x = 0
print(psi.rank(), pde.norm(psi, grid))                 # bond dimension, L2 norm
\end{lstlisting}

The units are the ones you put in. With $\hbar=m=1$ and $V=x^2/2$ the trap frequency
is $\omega=1$, so lengths are in units of the oscillator length
$a_{\rm ho}=\sqrt{\hbar/m\omega}$, time in $1/\omega$, energies in $\hbar\omega$, and
$g$ is the 1D coupling in those units. With $g=50$ the cloud is deep in the
Thomas--Fermi regime and the density approaches $n(x)=(\mu-x^2/2)/g$ with
$\mu=\tfrac12(3g/2)^{2/3}$. The script \code{examples/ground\_state\_gp.py} does that
comparison and gets $\mu$ within 1.5\,\% of the Thomas--Fermi value, the difference
being the kinetic (quantum-pressure) correction.

\code{solve} prints one line per step with the time, the bond dimensions and the
timing. Those bond dimensions are the single most useful diagnostic in the code, for
reasons that Section~\ref{sec:quantics} explains.

% ---------------------------------------------------------------------------
\section{Tensor networks in five minutes}
\label{sec:tn}

A \textbf{tensor} here is just an array with several indices, $T_{i_1 i_2\cdots i_n}$.
If every index takes two values the array has $2^n$ entries, and for $n=30$ that is a
billion numbers. A \textbf{tensor train} (in many-body physics, a \emph{matrix product
state}, MPS) writes such an array as a product of small matrices,
\begin{equation}
T_{i_1 i_2\cdots i_n} = A^{(1)}[i_1]\,A^{(2)}[i_2]\cdots A^{(n)}[i_n],
\label{eq:tt}
\end{equation}
where $A^{(s)}[i_s]$ is a $\chi_{s-1}\times\chi_s$ matrix for each value of $i_s$, with
$\chi_0=\chi_n=1$ so the product is a scalar. The matrix sizes $\chi_s$ are the
\textbf{bond dimensions}, or ranks. Storage is $\sum_s 2\chi_{s-1}\chi_s\sim 2n\chi^2$
instead of $2^n$: linear in $n$ rather than exponential, provided $\chi$ stays small.

Everything else follows from \eqref{eq:tt}.
\begin{itemize}[leftmargin=1.5em]
\item \textbf{Bond dimension measures the entanglement between the two halves of the
  index string.} Cut the train between sites $s$ and $s+1$; the array, reshaped as a
  matrix with the first $s$ indices as rows and the rest as columns, has at most rank
  $\chi_s$. A product array ($T=a_{i_1}b_{i_2}\cdots$) has $\chi=1$ everywhere. A
  generic random array has $\chi_s=2^{\min(s,\,n-s)}$ and gains nothing.
\item \textbf{Truncation is an SVD.} Any train can be brought to a canonical form in
  which the singular values of each cut are explicit. Discarding those below a
  tolerance, or beyond a maximum count \code{maxdim}, is the optimal low-rank
  approximation at that cut. This is how \code{tnde} keeps the field small after every
  operation: the discarded weight is the approximation error, and a train that keeps
  hitting \code{maxdim} is saying the field has more structure than was allowed for.
\item \textbf{Operators are trains too.} A matrix $O_{i_1\cdots i_n,\,j_1\cdots j_n}$
  written as a product of four-index cores $W^{(s)}[i_s,j_s]$ is a \textbf{matrix
  product operator} (MPO). Applying an MPO of rank $r$ to a train of rank $\chi$ gives,
  exactly, a train of rank $r\chi$, which is then truncated. A diagonal operator
  (multiplication by a function) is an MPO whose cores are diagonal in $(i_s,j_s)$.
\item \textbf{Contractions are cheap.} The overlap $\langle S|T\rangle$ of two trains,
  the norm, the sum of all entries and the value at a single index cost $O(n\chi^2)$ to
  $O(n\chi^3)$. This is how \code{norm}, \code{integral} and \code{evaluate} of
  Section~\ref{sec:out} work.
\end{itemize}

The last piece is how to \emph{build} a train for a function that can be evaluated but
not stored. \textbf{Tensor cross interpolation} (TCI) does that from samples. For a
matrix, choosing $\chi$ rows $I$ and $\chi$ columns $J$ gives the cross approximation
\begin{equation}
M \approx M_{:,J}\,(M_{I,J})^{-1}\,M_{I,:},
\end{equation}
exact if $M$ has rank $\chi$ and good if the chosen rows and columns (the
\textbf{pivots}) capture the dominant structure. Applied cut by cut along the train, a
few thousand function evaluations are enough to build a rank-$\chi$ train for a
function on $2^{30}$ points. TCI grows the pivot set greedily until the interpolation
error \emph{at the places it has looked} drops below a tolerance. That qualifier is the
caveat Section~\ref{sec:checklist} returns to.

% ---------------------------------------------------------------------------
\section{Quantics: a function as a tensor train}
\label{sec:quantics}

Take a function $f(x)$ on the interval $[x_{\min},x_{\min}+L)$ sampled on $M=2^R$
points, $x_j=x_{\min}+jL/2^R$. Write the index $j$ in binary,
\begin{equation}
j=\sum_{r=1}^{R} b_r\,2^{R-r},\qquad b_r\in\{0,1\},
\end{equation}
so that $b_1$ is the most significant bit and decides which half of the box $x$ is in,
$b_2$ which quarter, and so on down to $b_R$, which selects between neighbouring grid
points. The samples $f(x_j)$ are now an $R$-index tensor $F_{b_1b_2\cdots b_R}$, and
that tensor can be written as a train \eqref{eq:tt}. This is the \textbf{quantics}
representation: each site of the train is one \emph{length scale} of the function,
from the coarsest at site 1 to the finest at site $R$. Doubling the resolution adds one
site, not a factor of two in storage.

Which functions are low rank? Those whose behaviour at fine scales is simple.

\begin{tabularx}{\textwidth}{@{}lcL@{}}
\toprule
function & bond dimension & why \\
\midrule
$e^{ikx}$ & 1 & $e^{ikx_j}=\prod_r e^{ik\,b_r 2^{R-r}\Delta x}$ is a product over bits \\
polynomial of degree $p$ & $p+1$ & sums of products of bit-polynomials \\
$e^{-x^2/2\sigma^2}$ & $\approx 5$--$12$ at $10^{-10}$ & a Gaussian is a superposition of a few plane waves at every scale \\
a wave packet in a trap plus a lattice & 10--20 & a few smooth envelopes times a few oscillations \\
a smooth function with $n$ separate features & grows with $n$ & each feature costs its own rank at the fine scales \\
a step, $|x|$, a sharp front in 1D & small & a discontinuity on a line costs little \\
a curved sharp front in 2D & 50--100 & the interface position is a different fine-scale function along every row \\
noise, turbulence, chaos & $\sim 2^{R/2}$ & nothing to compress \\
\bottomrule
\end{tabularx}

\begin{finding}
\textbf{The quantics rank is set by the structure of the field, not by how finely it
is resolved.} A Gaussian on $2^{12}$ or on $2^{30}$ points has the same rank. Refining
the grid is free; adding physics that puts independent information at every length
scale is not.
\end{finding}

\paragraph{Several dimensions.} For $u(x,y)$ on $2^R\times2^R$ points the bits of both
variables are \textbf{interleaved} by length scale,
\begin{equation}
b^x_1\,b^y_1\;\;b^x_2\,b^y_2\;\;\cdots\;\;b^x_R\,b^y_R,
\end{equation}
so the train has $dR$ sites and the two bits at any given scale are neighbours. That
ordering keeps a smooth $d$-dimensional function compact. It has a consequence visible
in the logs: a separable field $f(x)g(y)$ has rank $r_f r_g$ at the bonds between
scales, so 2D and 3D ranks of 40--70 are normal for fields that would be rank 8 in 1D.
The solver switches to a cheaper, variational way of applying operators in 2D and 3D
for exactly this reason.

\paragraph{Momentum.} The grid is periodic with period exactly $L$, so the discrete
Fourier transform of the samples is the set of plane waves $e^{ik_jx}$ with
\begin{equation}
k_j=\frac{2\pi j}{L}\ \ (j<2^{R-1}),\qquad k_j=\frac{2\pi (j-2^R)}{L}\ \ (j\ge2^{R-1}),
\end{equation}
the ordinary physical momentum, and the largest representable momentum is
$\kmax=\pi/\Delta x=\pi 2^R/L$. The Fourier transform of a quantics train is itself an
MPO of small rank (about 12 at $10^{-12}$), an analytic construction rather than a fit,
so switching between position and momentum costs one MPO application. A
translation-invariant operator $\ell(-i\nabla)$ is applied exactly by transforming,
multiplying by the diagonal MPO of $\ell(\vk)$, and transforming back. The momentum grid
is the same interleaved train as the position grid, with the variables in reverse
order (a detail \code{Grid} handles).

\code{Grid} is the object that owns all of this: index $\leftrightarrow$ bits, index
$\leftrightarrow$ coordinate, index $\leftrightarrow$ momentum, the spacing and the
volume element.
\begin{lstlisting}
grid = Grid(12, [(-20.0, 20.0), (-20.0, 20.0)])
grid.M, grid.ndim, grid.nsites   # 4096, 2, 24
grid.spacings, grid.dvol         # [0.0098, 0.0098], 9.5e-5
x, y = grid.mesh()               # dense coordinate arrays: only for small R
kx, ky = grid.kmesh()            # dense momentum arrays
\end{lstlisting}
Indices are 64-bit integers, so a grid has at most 62 sites: $R\le31$ per axis in 1D
and 2D, $R\le20$ in 3D.

% ---------------------------------------------------------------------------
\section{How a time step works}
\label{sec:step}

The three terms of \eqref{eq:master} are applied in turn by \textbf{Strang splitting}.
Over one step of length $h$,
\begin{equation}
u(t+h)\;\approx\;N_{h/2}\;P_{h/2}\;L_{h}\;P_{h/2}\;N_{h/2}\;u(t),
\label{eq:strang}
\end{equation}
where $L_h=e^{h\ell(-i\nabla)}$, $P_h=e^{hp(\vx)}$, and $N_h$ is the exact flow of the
local ODE $\dot u=N(u,\vx,t)$ over time $h$. The symmetric arrangement makes the local
error $O(h^3)$ and the global error $O(h^2)$, regardless of what the three operators
are. The half-steps of $N$ at the end of one step and the start of the next are merged
into one full step wherever nothing needs the true state at the step boundary, so a
long run costs one nonlinear application per step; a snapshot or a renormalisation pays
the extra half step where it occurs.

Each factor is applied differently, and this is the part worth understanding.
\begin{itemize}[leftmargin=1.5em]
\item \textbf{$L_h$ is exact.} Fourier MPO, multiply by the diagonal MPO of
  $e^{h\ell(\vk)}$, inverse Fourier MPO. The only approximation is the truncation of the
  result to \code{maxdim} at \code{tolerance}. The multiplier MPO is cross-interpolated
  once, at start-up.
\item \textbf{$P_{h/2}$ is exact.} The diagonal MPO of $e^{hp(\vx)/2}$,
  cross-interpolated once. Again the only approximation is truncation.
\item \textbf{$N_h$ is applied by re-interpolation.} There is no fixed operator for a
  nonlinear map, so at every step the solver hands TCI the function ``evaluate the
  current train at this point, advance the local ODE there by $h$'' and lets it build a
  fresh train. The local ODE is advanced by a \textbf{closed-form flow map} when the
  equation provides one (every preset that has one does) and by classical \textbf{RK4}
  otherwise. This is a full cross interpolation from scratch, every step, and it
  dominates the cost of every nonlinear run.
\end{itemize}

Two physics-facing consequences follow.

\paragraph{Imaginary time.} With $t=-i\tau$ the Schr\"odinger equation becomes
$\partial_\tau\psi=-H\psi$, whose solution $e^{-\tau H}\psi_0$ projects any initial
state onto the ground state if it is renormalised as it decays.
\code{imaginary\_time=True} in the Schr\"odinger-type presets makes the symbol real
($-k^2/2m$), the potential real ($-V$), the nonlinear flow the exact decay of
$|\psi|^2$, and switches on renormalisation after every step. For a mixture,
\code{normalise="each"} holds every component's population fixed instead. The converged
state satisfies the (Gross--Pitaevskii) eigenvalue problem up to the $O(\Delta\tau^2)$
splitting error, which is why the harmonic-oscillator test in the \code{README} lands
at $10^{-4}$ and not $10^{-12}$: reduce $\Delta\tau$ to reduce it.

\paragraph{The momentum window.} A diffusive multiplier $e^{-hDk^2}$ decays to zero at
large $k$ and is a low-rank function of the momentum bits. An oscillatory one,
$e^{-ihk^2/2m}$, is not: at the largest momentum of a fine grid its phase changes by
many radians from one index to the next, and no small train represents that. The
Schr\"odinger-type presets therefore take a cutoff \code{kcut}, applied as a separable
Fermi--Dirac window
\begin{equation}
w(\vk)=\prod_v\frac{1}{1+e^{\beta(|k_v|-\kcut)}},
\qquad \beta=20/\kcut\ \text{by default},
\end{equation}
multiplied into the propagator.
\begin{finding}
\textbf{The window discards all momentum content above $\kcut$.} It is a rank-control
device with a physical cost, and it is not norm-preserving. Choose $\kcut$ above every
momentum the dynamics can excite (the lattice momentum, $\sqrt{2m(E-V)}$ for the
fastest part of a packet, $\sim\sqrt{2m\mu}$ for the healing length of a condensate),
check the norm \emph{without} renormalisation, and move the cutoff to see what
changes. In the published 1D benchmark the window, not the bond dimension, was the
leading approximation (\code{docs/guide.tex}). Diffusive equations need no window.
\end{finding}

The dense solver \code{tnde.reference} runs the \emph{identical} splitting loop and
local flow on full arrays with \code{numpy.fft}. Every difference between it and the
train solver is truncation. Use it on a small grid to separate ``the tensor network is
losing something'' from ``the time step is too big''.

% ---------------------------------------------------------------------------
\section{Setting up a problem}
\label{sec:setup}

\subsection{The grid}
\begin{lstlisting}
Grid(R, (lo, hi))                        # 1D, 2**R points on [lo, hi)
Grid(R, [(lo, hi), (lo2, hi2)])          # 2D, 2**R per axis
Grid(R, [(lo, hi)] * 3)                  # 3D
\end{lstlisting}
The box is periodic and the upper endpoint is excluded. Two physical constraints set
$L$ and $R$: the box must be large enough that the field and its periodic images do not
interact (a decaying field should be negligible at the edges), and the spacing
$\Delta x=L/2^R$ must resolve the finest physical length (a lattice period, a healing
length $\xi=1/\sqrt{2mgn}$, a diffusion front $\sqrt{Dt}$). Because refining costs one
site per doubling, take $R$ generous: $R=20$ in 1D is $10^6$ points and costs the same
as $R=12$ for a smooth field.

\subsection{The equation}

A preset, or the three callables directly. The presets in \code{tnde.equations}:

\begin{small}
\begin{tabularx}{\textwidth}{@{}p{1em}LL@{}}
\toprule
\multicolumn{3}{@{}l}{preset} \\
 & equation & local step \\
\midrule
\multicolumn{3}{@{}l}{\code{schrodinger(V, m, hbar, kcut)}} \\
 & $i\hbar\partial_t\psi=-\frac{\hbar^2}{2m}\nabla^2\psi+V\psi$ & none \\[3pt]
\multicolumn{3}{@{}l}{\code{gross\_pitaevskii(V, g, m, kcut, imaginary\_time)}} \\
 & $i\partial_t\psi=-\frac{1}{2m}\nabla^2\psi+V\psi+g|\psi|^2\psi$ & exact phase $e^{-ig|\psi|^2h}$; in imaginary time the exact decay $\psi/\sqrt{1+2g|\psi|^2h}$ \\[3pt]
\multicolumn{3}{@{}l}{\code{nonlinear\_schrodinger(f, V, m, kcut)}} \\
 & $i\partial_t\psi=-\frac{1}{2m}\nabla^2\psi+V\psi+f(|\psi|^2)\psi$ & exact phase $e^{-if(|\psi|^2)h}$ \\[3pt]
\multicolumn{3}{@{}l}{\code{coupled\_gross\_pitaevskii(G, V, m, kcut, imaginary\_time)}} \\
 & $i\partial_t\psi_a=-\frac{1}{2m}\nabla^2\psi_a+V_a\psi_a+\sum_bG_{ab}|\psi_b|^2\psi_a$ & exact phases; RK4 in imaginary time \\[3pt]
\multicolumn{3}{@{}l}{\code{heat(D, source, rate)}} \\
 & $\partial_tu=D\nabla^2u+r(\vx)u+S(u,\vx,t)$ & RK4 if a source is given \\[3pt]
\multicolumn{3}{@{}l}{\code{fractional\_diffusion(alpha, D)}} \\
 & $\partial_tu=-D(-\nabla^2)^{\alpha/2}u$ & none \\[3pt]
\multicolumn{3}{@{}l}{\code{advection\_diffusion(velocity, D)}} \\
 & $\partial_tu=-\mathbf c\cdot\nabla u+D\nabla^2u$ & none; advection is the exact multiplier $-i\mathbf c\cdot\vk$ \\[3pt]
\multicolumn{3}{@{}l}{\code{fisher\_kpp(D, r)}} \\
 & $\partial_tu=D\nabla^2u+ru(1-u)$ & exact logistic map $u/(u+(1-u)e^{-rh})$ \\[3pt]
\multicolumn{3}{@{}l}{\code{allen\_cahn(eps)}} \\
 & $\partial_tu=\varepsilon\nabla^2u+u-u^3$ & exact $u/\sqrt{u^2+(1-u^2)e^{-2h}}$ \\[3pt]
\multicolumn{3}{@{}l}{\code{complex\_ginzburg\_landau(b, c)}} \\
 & $\partial_tA=A+(1+ib)\nabla^2A-(1+ic)|A|^2A$ & exact $A(1+2|A|^2h)^{-(1+ic)/2}$ \\[3pt]
\multicolumn{3}{@{}l}{\code{gray\_scott(Du, Dv, F, k)}} \\
 & $\partial_tu=D_u\nabla^2u-uv^2+F(1-u)$, $\partial_tv=D_v\nabla^2v+uv^2-(F+k)v$ & RK4 \\[3pt]
\multicolumn{3}{@{}l}{\code{fitzhugh\_nagumo(Du, Dv, a, eps, gamma)}} \\
 & $\partial_tu=D_u\nabla^2u+u(1-u)(u-a)-v$, $\partial_tv=D_v\nabla^2v+\varepsilon(u-\gamma v)$ & RK4 \\[3pt]
\bottomrule
\end{tabularx}
\end{small}

Every symbol is written for any number of dimensions, so the same preset serves 1D, 2D
and 3D. Potentials and sources are functions of the coordinates, \code{V(x)} in 1D,
\code{V(x, y)} in 2D, evaluated on NumPy arrays.

A custom equation supplies the three maps. The conventions: \code{linear} receives the
momentum components and returns $\partial_t\hat u_{\vk}/\hat u_{\vk}$; \code{potential}
receives the coordinates and returns $\partial_tu/u$; \code{nonlinear} receives the
field, the coordinates and the time, in that order, and returns $\partial_tu$.
\begin{lstlisting}
from tnde import Equation

# 1D:  d_t u = -0.01 k^4 u  -  x^2 u  +  i a x cos(w t) u
eq = Equation(linear=lambda k: -0.01 * k**4,
              potential=lambda x: -x**2,
              nonlinear=lambda u, x, t: 1j * a * x * np.cos(w * t) * u)
\end{lstlisting}
Time-dependent local terms belong in \code{nonlinear}; \code{potential} is static so
that its MPO is built once. If the exact solution of the local ODE is known, pass it as
\code{flow(u, x, t, h)} instead of \code{nonlinear}: one evaluation per point instead
of four per RK4 sub-step, and no time-stepping error. The presets show the pattern. If
RK4 is used and the local dynamics is stiff, raise \code{rk\_substeps}.

A multi-component field sets \code{ncomp}, passes and returns tuples in the local
terms, and takes a symbol per component (or one shared symbol):
\begin{lstlisting}
# two species, 2D, different diffusivities, a coupled local ODE
eq = Equation(linear=[lambda kx, ky: -Du * (kx**2 + ky**2),
                      lambda kx, ky: -Dv * (kx**2 + ky**2)],
              nonlinear=lambda uv, x, y, t: (f(*uv), g(*uv)),
              ncomp=2)
\end{lstlisting}
The remaining fields are \code{filter} (a momentum window $w(\vk)$, which the
Schr\"odinger presets set from \code{kcut}), \code{normalise} (\code{True} for unit
total norm after every step, \code{"each"} for unit norm per component) and
\code{rk\_substeps}.

\subsection{The initial condition}

A callable of the coordinates, or a list of them for a multi-component field. It is
cross-interpolated onto the grid at \code{tolerance} and \code{maxdim} before the first
step, and the log reports its rank. An existing train, for instance the result of a
previous run, can be passed instead.

\subsection{Running}
\begin{lstlisting}
u, snaps = pde.solve(eq, u0, grid, dt, nsteps,
                     tolerance=1e-10,   # truncation and TCI tolerance of the field
                     maxdim=32,         # bond-dimension cap of the field
                     save_every=None,   # keep every k-th state in `snaps`
                     callback=None,     # callback(step, t, state) at every snapshot
                     method="auto",     # how diagonal MPOs are applied
                     t0=0.0, verbose=True)
\end{lstlisting}
\code{u} is the final state, a \code{TensorTrain} (or a list of them). \code{snaps} is
a dictionary keyed by step number, with the initial state under 0 and the state at
$t_0+j\,\dd t$ under $j$. Snapshots are true states at step boundaries, so a run with
snapshots costs one extra half nonlinear step per snapshot. A callback is the cheap
way to monitor an observable during a long run:
\begin{lstlisting}
def report(step, t, u):
    (x,), v = pde.reconstruct(u, grid, n=512)
    print(f"t = {t:6.2f}  rank {u.rank():3d}  peak density {np.max(np.abs(v)**2):.4f}")

u, _ = pde.solve(eq, u0, grid, dt, nsteps, save_every=20, callback=report, verbose=False)
\end{lstlisting}

% ---------------------------------------------------------------------------
\section{Getting physics out}
\label{sec:out}

The state is a train, so nothing is a plain array until asked for. All of these live in
\code{tnde.pde}.

\begin{tabularx}{\textwidth}{@{}lLl@{}}
\toprule
call & returns & cost \\
\midrule
\code{reconstruct(u, grid, n, window)} & \code{(axes, values)}: the field on $n$ points
per direction across the box, or across a \code{window} of \code{(lo, hi)} pairs &
$n^d$ evaluations \\
\code{evaluate(u, grid, *x)} & the field at the grid points nearest to real coordinates,
broadcast over arrays & one contraction per point \\
\code{to\_dense(u, grid)} & the full $(2^R)^d$ array & small grids only \\
\code{norm(u, grid)} & $\sqrt{\sum_c\int|u_c|^2\,\dd V}$ & $O(R\chi^3)$ \\
\code{integral(u, grid)} & $\int u\,\dd V$ & $O(R\chi^2)$ \\
\code{normalise(u, grid, each)} & the state rescaled to unit norm & \\
\code{interpolate(fn, grid, tolerance, maxdim)} & a train of an arbitrary function & one TCI \\
\code{max\_error(u, fn, grid, n, window)} & maximum relative deviation of \code{u} from
\code{fn} on an independent sample & $n^d$ evaluations \\
\code{fourier\_transform(u, grid)} & the train in momentum space & one MPO application \\
\bottomrule
\end{tabularx}

Densities and plots come from \code{reconstruct}. It samples, so choose $n$ to resolve
what is to be seen and use \code{window} to zoom:
\begin{lstlisting}
(x, y), psi_xy = pde.reconstruct(psi, grid, n=512, window=[(-5, 5), (-5, 5)])
density = np.abs(psi_xy)**2
\end{lstlisting}

\paragraph{Expectation values.} The exact route stays on the train: multiply the
density by the observable and integrate, with the density evaluated on grid points so
that nothing is interpolated between them.
\begin{lstlisting}
def expectation(psi, grid, obs):
    """<obs> = int obs(x) |psi|^2 dV / int |psi|^2 dV, for a diagonal observable."""
    dens = lambda *x: obs(*x) * np.abs(pde.evaluate(psi, grid, *x))**2
    return (pde.integral(pde.interpolate(dens, grid, tolerance=1e-10), grid)
            / pde.norm(psi, grid)**2).real

x_mean  = expectation(psi, grid, lambda x: x)
x2_mean = expectation(psi, grid, lambda x: x**2)
width   = np.sqrt(x2_mean - x_mean**2)
\end{lstlisting}
Momentum-space quantities go through \code{fourier\_transform} and the same integrals
in $k$; the kinetic energy of the first run is $\frac{1}{2m}\int k^2|\hat\psi_k|^2$. On
a small grid it is simpler, and exact, to take \code{to\_dense} and use NumPy directly,
which is what the examples do for their reference numbers.

\paragraph{Checking against a dense solution.} For $R\lesssim16$ in 1D or
$R\lesssim8$ in 2D the whole problem fits in memory, and
\begin{lstlisting}
from tnde import reference
u_ref, _ = reference.solve(eq, u0, grid, dt, nsteps)
err = np.max(np.abs(pde.to_dense(u, grid) - u_ref)) / np.max(np.abs(u_ref))
\end{lstlisting}
isolates truncation error from everything else. Do this once, at small $R$, before
scaling up. The ranks do not depend on $R$ for a smooth field, so what is learnt at
$R=14$ transfers to $R=30$.

% ---------------------------------------------------------------------------
\section{Choosing the parameters}
\label{sec:params}

\paragraph{\code{dt}.} The splitting error is $O(\dd t^2)$ globally and independent of
the grid: there is no CFL condition, because $L_h$ is exact for any $h$. What limits
\code{dt} is the physics: the phase $\omega\,\dd t$ of the fastest relevant oscillation
should be small, and in imaginary time $\dd\tau\,(E_{\max}-E_0)\lesssim1$ for the
highest energy present. The quick convergence test is to halve \code{dt}; the error
should drop by four.

\paragraph{\code{tolerance}.} Both the SVD truncation of the field and the TCI tolerance
of the nonlinear step. Values of $10^{-8}$ to $10^{-10}$ are normal. Every step re-fits
the field, so errors accumulate linearly in the number of steps; the \code{README}
tables show $10^{-12}$ agreement with dense references over tens of steps at
$10^{-10}$.

\paragraph{\code{maxdim}.} The bond-dimension cap. It caps cost, which scales as $\chi^2$
to $\chi^3$ per step, and it caps accuracy: once a rank hits \code{maxdim} the
truncation is silent. Watch the ranks in the log. If they sit well below the cap, the
cap is inactive; if they sit at it, either raise it or accept that the result is a
low-rank projection of the dynamics. Typical: 20--30 in 1D, 40--70 in 2D and 3D.

\paragraph{\code{method}.} How the diagonal MPOs ($P_h$ and the multiplier) are applied.
\code{'naive'} forms the exact rank-$r\chi$ product and truncates it: cheap and slightly
more accurate when both ranks are small. \code{'fit'} sweeps a variational ansatz of
bond dimension \code{maxdim}: about $50\times$ faster when the ranks are 50--70.
\code{'auto'} picks \code{'naive'} in 1D and \code{'fit'} in 2D and 3D. The Fourier MPOs
always use the exact product.

\paragraph{\code{kcut}.} See the momentum window in Section~\ref{sec:step}. Only
Schr\"odinger-type equations need it, and only on grids fine enough that the free
phase per step at $\kmax$, $h\,\kmax^2/2m$ with $\kmax=\pi2^R/L$, exceeds a few radians.
Place it above every momentum the physics excites, with margin, and verify by checking
the norm without renormalisation and by moving it.

\paragraph{Cost.} One core, single-threaded BLAS, from the committed examples:

\begin{tabularx}{\textwidth}{@{}Lll@{}}
\toprule
case & grid & per step \\
\midrule
1D Bose gas ground state, $g=50$ & $2^{16}$ & 0.17 s \\
two-component mixture, RK4 local step & $2^{14}$ & 0.36 s \\
2D advection--diffusion of a Gaussian, rank 40 & $4096^2$ & 2.1 s \\
2D Allen--Cahn droplet, rank 50 & $128^2$ & 1.0 s \\
1D Gross--Pitaevskii in a lattice (benchmark) & $2^{30}$ & 3 s \\
2D Gross--Pitaevskii in a quasicrystal (benchmark) & $2^{20}\times2^{20}$ & 73 s \\
\bottomrule
\end{tabularx}

Neither column is the number of grid points. The nonlinear re-interpolation costs
(number of TCI evaluations) $\times$ (rank$^2$) $\times$ (cost of one local flow), and
the ranks are set by the structure of the field. The Allen--Cahn droplet on $128^2$
points takes as long as the Gaussian on $4096^2$ because its curved interface is the
case the representation is not for.

% ---------------------------------------------------------------------------
\section{Worked examples and the physics in them}
\label{sec:examples}

All in \code{examples/}, each checked against a dense reference or an exact solution.
\begin{itemize}[leftmargin=1.5em]
\item \code{ground\_state\_gp.py} relaxes a 1D Bose gas in a harmonic trap to its
  ground state in imaginary time and compares with the Thomas--Fermi profile. The
  residual difference in $\mu$ is the quantum-pressure term
  $-\tfrac12(\sqrt n)''/\sqrt n$, which the Thomas--Fermi approximation drops.
\item \code{binary\_mixture\_1d.py} is a two-component condensate with
  $G=g\left(\begin{smallmatrix}1&r\\r&1\end{smallmatrix}\right)$, $r>1$, relaxed with
  each population held fixed. Above the miscibility threshold $G_{12}^2>G_{11}G_{22}$
  the species expel each other and settle side by side; the reported overlap integral
  $\int n_1n_2\,\dd x$ drops from 0.26 to 0.02 and the centres of mass move to $\pm1.5$.
\item \code{advection\_diffusion\_2d.py} is the case the representation is built for:
  a Gaussian drifting with velocity $\mathbf c$ and spreading as $\sigma^2+2Dt$ on a
  $4096^2$ grid, compared with the closed form at every snapshot to $2\times10^{-7}$ at
  rank 40. Run it with $R=15$ to see $10^9$ points evolve on a laptop.
\item \code{allen\_cahn\_droplet\_2d.py} is the opposite case, kept to show the cost: a
  circular droplet of the $u=+1$ phase shrinking by curvature flow,
  $r^2=r_0^2-2\varepsilon t$, reproduced to 0.5\,\%. The interface is sharp and curved,
  so the rank is 50 at $10^{-5}$ and 90 at $10^{-8}$ on a $128^2$ grid that NumPy would
  finish in milliseconds.
\item \code{reproduce\_paper\_1d.py}, \code{reproduce\_paper\_2d.py} reproduce the
  published quantics Gross--Pitaevskii benchmarks (Section~\ref{sec:legacy}).
\end{itemize}
A convenient pattern for a new problem is to copy the nearest example, drop $R$ to
something dense-checkable, confirm the reference agreement, then raise $R$.

% ---------------------------------------------------------------------------
\section{Checklist before trusting a number}
\label{sec:checklist}

\begin{enumerate}[leftmargin=1.5em]
\item \textbf{Ranks below the cap?} A rank pinned at \code{maxdim} means silent
  truncation. Raise the cap or lower the tolerance until the ranks stop changing.
\item \textbf{Converged in \code{dt}?} Halve it once. The splitting error is
  $O(\dd t^2)$.
\item \textbf{Box large enough, grid fine enough?} Double $L$ (with $R+1$, to keep
  $\Delta x$) and check nothing changes. The periodic boundary is invisible only if the
  field vanishes at the edges.
\item \textbf{Window above the physics?} With \code{kcut} set, run once with
  \code{normalise} off and look at the norm. A drift means momentum is being thrown
  away. A renormalised norm of 1.000000 proves nothing.
\item \textbf{TCI looked everywhere?} TCI's error estimate is taken at the pivots it
  chose. A feature no pivot reached is invisible to it. The solver seeds both sides of
  the box centre (a centred feature sits exactly on the most significant bit of every
  variable) and the peaks of $|f|$, and runs a global pivot search after every sweep,
  but the honest check is \code{max\_error} against an independent sample, or a dense
  reference at small $R$.
\item \textbf{Dense reference agreed at small $R$?} The ranks are $R$-independent for a
  smooth field, so agreement at $R=14$ is evidence at $R=30$.
\item \textbf{Real equation, complex field.} Everything is \code{complex128}. Take
  \code{.real} of what is reconstructed.
\end{enumerate}

% ---------------------------------------------------------------------------
\section{The Gross--Pitaevskii benchmark drivers}
\label{sec:legacy}

The library grew out of a solver for a published quantics tensor-train
Gross--Pitaevskii benchmark (\href{https://arxiv.org/abs/2507.04262}{arXiv:2507.04262}),
and those drivers survive as \code{tnde.evolve1d}, \code{tnde.evolve2d}, their dense
oracle \code{tnde.dense} and \code{tnde.observables}. They reproduce the published
tensors: at $R=30$ on $[-500,500]$, a Gaussian in a harmonic trap plus a fast lattice,
$g=5$, $\Delta t=0.01$, the bond dimensions match step for step and the overlap with
the reference state is $1-2\times10^{-10}$ after 100 steps.

They keep the benchmark's conventions, which differ from \code{pde}: a grid that
includes the endpoint, the finite-difference dispersion
$-\frac{4M^2}{L^2}\sin^2(\pi k/M)$ instead of $-k^2$, a renormalisation after every
kinetic step, and a momentum cutoff searched on a Gaussian trial state. Use them to
reproduce that benchmark; use \code{pde} for everything else. Their physics guide is
\code{docs/guide.tex}, whose error budget is worth reading before any Schr\"odinger-type
run: it finds that the momentum window, not the bond dimension, is the leading
approximation in the 1D benchmark, with the cutoff a factor of six below the momentum
the lattice drives and the renormalisation hiding the loss.

% ---------------------------------------------------------------------------
\section{Symbols}

\begin{tabularx}{\textwidth}{@{}lL@{}}
\toprule
symbol & meaning \\
\midrule
$R$ & bits per space dimension; $2^R$ grid points per axis \\
$M=2^R$ & grid points per axis \\
$L$, $\Delta x=L/2^R$ & box length and spacing; the box is periodic and excludes its upper end \\
$b_r$ & the $r$-th binary digit of a grid index; site $r$ of the train is length scale $L/2^r$ \\
$\chi$ & bond dimension (rank) of a train; \code{maxdim} is its cap \\
$\ell(\vk)$ & Fourier symbol of the linear operator, \code{linear} \\
$p(\vx)$ & local linear coefficient, \code{potential} \\
$N(u,\vx,t)$ & local nonlinearity, \code{nonlinear}; its exact flow is \code{flow} \\
$h$, \code{dt} & time step; splitting error $O(h^2)$ \\
$\kmax=\pi/\Delta x$ & largest representable momentum \\
$\kcut$, $\beta$ & Fermi--Dirac momentum window position and inverse width \\
$\tau=it$ & imaginary time; relaxation to the ground state \\
TCI & tensor cross interpolation: builds a train from function samples at pivots \\
MPO & matrix product operator: an operator written as a train \\
\bottomrule
\end{tabularx}

\end{document}
